\documentclass[10pt,a4paper]{amsart}
\usepackage[margin=19mm]{geometry}
\usepackage{amsmath,amssymb,mathtools}
\usepackage[scr=rsfs,cal=euler]{mathalfa}
\usepackage{xfrac}
\usepackage[unicode,hidelinks,bookmarks=false]{hyperref}
\newcommand{\ie}{i.e.\ }
\newcommand{\cf}{cf.\ }
\newcommand{\resp}{respectively\ }
\newcommand{\Subsection}{\subsection}
\providecommand{\nobreakdash}{\nobreak\hbox{-}}
\setlength{\emergencystretch}{2em}
\multlinegap0pt
\usepackage{bm}
\usepackage{bbm}
\usepackage{dsfont}
\usepackage{xspace}
\usepackage{cancel}
\usepackage{mathtools}
\usepackage{amsthm}
\makeatletter
\@ifundefined{theorem}{\newtheorem{theorem}{Theorem}}{}
\@ifundefined{lemma}{\newtheorem{lemma}[theorem]{Lemma}}{}
\makeatother
\newcommand\OO{\mathcal{O}}
\newcommand\PP{\mathbb{P}}
\newcommand\QQ{\mathbb{Q}}
\title[Non-liftable varieties via etale cohomology rings]{Non-liftable varieties via etale cohomology rings}
\author{Runjie Hu, Siqing Zhang}
\date{Source: arXiv:2607.18588v1 (CC BY 4.0)}
\begin{document}
\maketitle

\noindent\emph{Source and licence.} Non-liftable varieties via etale cohomology rings. Version arXiv:2607.18588v1 (CC BY 4.0), distributed under the Creative Commons Attribution 4.0 International licence, as recorded in the arXiv licence metadata for this exact version and shown on the abstract page https://arxiv.org/abs/2607.18588v1. Full rights notice: licenses/arxiv-2607.18588-CC-BY-4.0.txt.\par\smallskip

\noindent\textit{Editorial scope.} Counted unit: the Lemma whose source label is \texttt{lem:distinguished-lines} in the article (source0.tex lines 971--996, source label \texttt{lem:distinguished-lines}), delivered together with the article's own complete original proof of it (source0.tex lines 997--1069, one \texttt{proof} environment of 73 lines). No mathematics is written, repaired or re-derived: statement and proof are the article's own text line for line. Required context retained uncounted: eqn:blowup-for-X (lines 862--870); eqn:ambient-exceptional-product (lines 895--900); eqn:exceptional-product (lines 902--907). Every definition, display and label of those lines is retained unchanged. Omitted from this package and not redistributed: 1--861, 871--894, 901--901, 908--970, 1070--1205. Nothing inside a retained excerpt is omitted or altered: every retained body is delivered exactly as the article's lines are written, so no omission receipt of any kind accompanies this package. Citations: the retained excerpts carry no citation command at all, so no bibliography entry of the article is transcribed and no reference is redistributed. Reference adaptation: none was needed and none was made; every reference inside the delivered unit resolves inside it, so no unresolved reference or citation is produced. Printed slips kept literal: no printed word, symbol, hypothesis or number of the counted statement or of its proof was altered. Layout and class: the article's own class is \texttt{amsart} with packages including amssymb, amsfonts, amsmath, graphicx, xcolor, mathrsfs, stmaryrd, epsfig, color, blindtext, enumerate, hyperref, url, bbm; the wrapper is the lane's pinned \texttt{amsart} editorial wrapper at 10pt on A4, which restates the theorem-like environments the retained text needs and adds the article's own macro definitions that the retained text needs (\texttt{\textbackslash OO}, \texttt{\textbackslash PP}, \texttt{\textbackslash QQ}), with extra wrapper packages (bm, bbm, dsfont, xspace, cancel, mathtools). The delivered numbering is the unit's own. Assets: no retained span contains a figure environment, a graphics-inclusion command, a PDF-page inclusion, a table or a TikZ picture; no image file is copied into this package, the package declares no asset dependency and no image file is redistributed with it. Licence: version arXiv:2607.18588v1 of this article is distributed under the Creative Commons Attribution 4.0 International licence, as recorded in the arXiv licence metadata for this exact version and shown on the abstract page https://arxiv.org/abs/2607.18588v1. A full rights notice is delivered at licenses/arxiv-2607.18588-CC-BY-4.0.txt. Characters: every retained excerpt is pure ASCII, so the delivered engine, XeLaTeX with its default Latin Modern text font, reports no missing character.
\begin{equation}
\label{eqn:blowup-for-X}
H^m(X)
=
\pi^*H^m(\PP^9_k)
\oplus
\bigoplus_{a=1}^{5}
\iota_a\bigl(H^{m-2a}(Z)\bigr).
\end{equation}

\begin{equation}
\label{eqn:ambient-exceptional-product}
\pi^*\delta\cdot\iota_a(\alpha)
=
\iota_a\bigl(i^*\delta\cdot\alpha\bigr).
\end{equation}

\begin{equation}
\label{eqn:exceptional-product}
\iota_a(\alpha)\cdot\iota_b(\beta)
=
\iota_{a+b}(\alpha\cdot\beta).
\end{equation}

\begin{lemma}[The two distinguished lines]
\label{lem:distinguished-lines}
The lines $\QQ_\ell h$ and $\QQ_\ell e$ are determined intrinsically by
the graded algebra $H^*(X)$.  More precisely,
\begin{equation}
\label{eqn:hyperplane-line}
\QQ_\ell h
=
\left\{
x\in H^2(X):x^3\cdot H^3(X)=0
\right\}.
\end{equation}
Once this line has been recovered, for any nonzero
$\widetilde h\in\QQ_\ell h$,
\begin{equation}
\label{eqn:exceptional-line}
\QQ_\ell e
=
\ker\left(
\widetilde h^{\,4}\cdot-:
H^2(X)\longrightarrow H^{10}(X)
\right).
\end{equation}
Both descriptions are compatible with extension of the coefficient
field.
\end{lemma}

\begin{proof}
By \eqref{eqn:blowup-for-X}, every class in $H^2(X)$ can be written
uniquely as $ah+be$.  We first show that the line $\QQ_\ell h$ is
contained in the right-hand side of \eqref{eqn:hyperplane-line}.  For
$u\in H^1(Z)$, equation \eqref{eqn:ambient-exceptional-product} gives
\[
h^3\cdot\iota_1(u)
=
\iota_1\left(
i^*c_1\bigl(\OO_{\PP^9_k}(1)\bigr)^3\cdot u
\right)
=
0,
\]
because the class inside $\iota_1$ belongs to $H^7(Z)$ and $Z$ is a
threefold.  Since $H^3(X)=\iota_1(H^1(Z))$, it follows that
$h^3\cdot H^3(X)=0$.

Conversely, suppose that $x=ah+be$ with $b\neq0$, and choose
$0\neq u\in H^1(Z)$.  
Expanding $x^3=(ah+be)^3$ and using
\eqref{eqn:exceptional-product}, we obtain
\[
\begin{aligned}
x^3\cdot\iota_1(u)
&=
a^3h^3\cdot\iota_1(u)
+3a^2b\,h^2e\cdot\iota_1(u)
+3ab^2\,he^2\cdot\iota_1(u)
+b^3e^3\cdot\iota_1(u)
\\
&=
a^3h^3\cdot\iota_1(u)
+3a^2b\,h^2\cdot\iota_2(u)
+3ab^2\,h\cdot\iota_3(u)
+b^3\iota_4(u).
\end{aligned}
\]
By \eqref{eqn:ambient-exceptional-product}, multiplication by $h$
preserves the index of an exceptional summand.  Thus the first three
terms belong respectively to the first, second, and third exceptional
summands in \eqref{eqn:blowup-for-X}, whereas the last term belongs to
$\iota_4(H^1(Z))$.  
Since $\iota_4$ is injective,
$b^3\iota_4(u)\neq0$.  Thus $x^3$ does not annihilate $H^3(X)$, proving
\eqref{eqn:hyperplane-line}.

It remains to recover the exceptional line.  Multiplying
$\widetilde h$ by a nonzero scalar does not change the kernel in
\eqref{eqn:exceptional-line}, so we may take $\widetilde h=h$.  Equation
\eqref{eqn:ambient-exceptional-product} gives
\[
h^4e
=
\iota_1\left(
i^*c_1\bigl(\OO_{\PP^9_k}(1)\bigr)^4
\right)
=
0,
\]
because $H^8(Z)=0$.  On the other hand, $h^5\neq0$, since it is the
pullback of the fifth power of the hyperplane class on $\PP^9_k$.
Therefore, we have that $h^4(ah+be)=ah^5$, whose kernel is $\QQ_\ell e$.  This proves
\eqref{eqn:exceptional-line}.

The condition in \eqref{eqn:hyperplane-line} is given by homogeneous
polynomial equations determined by multiplication in $H^*(X)$, and its
reduced common zero locus is the line $\QQ_\ell h$.  Once this line is
known, \eqref{eqn:exceptional-line} describes $\QQ_\ell e$ as the
kernel of a linear map determined by multiplication.  These
constructions are intrinsic to the graded algebra and commute with
extension of the coefficient field.
\end{proof}

\end{document}
